% !PRINTED 181
\weberpage{181}
6. \gesperrt{Eine Modulfunktion, die durch die Substitution $\left(\omega, -\dfrac{1}{\omega}\right)$ ungeändert bleibt und durch $(\omega, \omega+1)$ den Faktor $e^{\mp\frac{2\pi i}{3}}$ annimmt, ist das Produkt von $\gamma_3(\omega)^{\pm 1}$ mit einer rationalen Funktion $j(\omega)$.} % [sic] Printed gamma_3 although the stated cubic-root transformation requires gamma_2; also printed \textquotedbl{}Funktion j(\omega)\textquotedbl{} without \textquotedbl{}von\textquotedbl{}.

7. \gesperrt{Eine Modulfunktion, die durch die Substitutionen $\left(\omega, -\dfrac{1}{\omega}\right)$ das Zeichen ändert und durch $(\omega, \omega+1)$ den Faktor $-e^{\mp\frac{2\pi i}{3}}$ annimmt, ist das Produkt von $\gamma_2(\omega)\gamma_3(\omega)^{\pm 1}$ mit einer rationalen Funktion von $j(\omega)$.}

Eingehender werden wir uns mit den Modulfunktionen im zweiten Teil beschäftigen.

\begin{center}
\textbf{\S\ 55. Darstellung der elliptischen Funktionen durch $v$ und $\varkappa^2$.}
\end{center}

Wenn wir das Umkehrproblem der elliptischen Integrale so wie in \S\,51 als die Aufgabe der Integration der \gesperrt{elliptischen Differentialgleichungen} fassen, so verlangt es die Darstellung der drei Funktionen sn\,$v$, cn\,$v$, dn\,$v$ durch die beiden unabhängigen Variablen $v$, $\varkappa^2$. Diese Aufgabe ist durch \S\,42 gelöst, aber bezüglich der Variablen $\varkappa^2$ nur implizite. Man kann aber auch Darstellungen finden, in denen $\varkappa^2$ explizite vorkommt, und zwar durch Reihen, die nach Potenzen von $v$ fortschreiten, deren Koeffizienten \gesperrt{rationale Funktionen von $\varkappa^2$ sind}. Die Koeffizienten dieser Reihen können freilich nicht durch ein allgemeines Gesetz dargestellt, sondern nur sukzessive berechnet werden. Diese Darstellungen verdankt man hauptsächlich \gesperrt{Weierstrass}\footnote{Über die \gesperrt{Weierstrasssche} Theorie der elliptischen Funktionen vgl. man besonders die von H.\,A. \gesperrt{Schwarz} herausgegebenen \glqq Formeln und Lehrsätze zum Gebrauch der elliptischen Funktionen\grqq. Über die Reihenentwickelungen vgl. man auch \gesperrt{Königsberger}, \glqq Vorlesungen über die Theorie der elliptischen Funktionen\grqq, 25. Vorlesung.}.

Die $\sigma$-Funktionen können, da sie durchaus endliche und stetige Funktionen von $u$ sind, in unbedingt konvergente Potenzreihen nach $u$ entwickelt werden, und wenn wir daher die in \S\,45, (2) vorkommenden $\sigma$-Funktionen in dieser Weise entwickeln, so bekommen wir eine Lösung unserer Aufgabe. Wir setzen daher
%CONT>>
